% Options for packages loaded elsewhere
\PassOptionsToPackage{unicode}{hyperref}
\PassOptionsToPackage{hyphens}{url}
%
\documentclass[
]{article}
\usepackage{amsmath,amssymb}
\usepackage{iftex}
\ifPDFTeX
  \usepackage[T1]{fontenc}
  \usepackage[utf8]{inputenc}
  \usepackage{textcomp} % provide euro and other symbols
\else % if luatex or xetex
  \usepackage{unicode-math} % this also loads fontspec
  \defaultfontfeatures{Scale=MatchLowercase}
  \defaultfontfeatures[\rmfamily]{Ligatures=TeX,Scale=1}
\fi
\usepackage{lmodern}
\ifPDFTeX\else
  % xetex/luatex font selection
\fi
% Use upquote if available, for straight quotes in verbatim environments
\IfFileExists{upquote.sty}{\usepackage{upquote}}{}
\IfFileExists{microtype.sty}{% use microtype if available
  \usepackage[]{microtype}
  \UseMicrotypeSet[protrusion]{basicmath} % disable protrusion for tt fonts
}{}
\makeatletter
\@ifundefined{KOMAClassName}{% if non-KOMA class
  \IfFileExists{parskip.sty}{%
    \usepackage{parskip}
  }{% else
    \setlength{\parindent}{0pt}
    \setlength{\parskip}{6pt plus 2pt minus 1pt}}
}{% if KOMA class
  \KOMAoptions{parskip=half}}
\makeatother
\usepackage{xcolor}
\usepackage[margin=1in]{geometry}
\usepackage{color}
\usepackage{fancyvrb}
\newcommand{\VerbBar}{|}
\newcommand{\VERB}{\Verb[commandchars=\\\{\}]}
\DefineVerbatimEnvironment{Highlighting}{Verbatim}{commandchars=\\\{\}}
% Add ',fontsize=\small' for more characters per line
\usepackage{framed}
\definecolor{shadecolor}{RGB}{248,248,248}
\newenvironment{Shaded}{\begin{snugshade}}{\end{snugshade}}
\newcommand{\AlertTok}[1]{\textcolor[rgb]{0.94,0.16,0.16}{#1}}
\newcommand{\AnnotationTok}[1]{\textcolor[rgb]{0.56,0.35,0.01}{\textbf{\textit{#1}}}}
\newcommand{\AttributeTok}[1]{\textcolor[rgb]{0.13,0.29,0.53}{#1}}
\newcommand{\BaseNTok}[1]{\textcolor[rgb]{0.00,0.00,0.81}{#1}}
\newcommand{\BuiltInTok}[1]{#1}
\newcommand{\CharTok}[1]{\textcolor[rgb]{0.31,0.60,0.02}{#1}}
\newcommand{\CommentTok}[1]{\textcolor[rgb]{0.56,0.35,0.01}{\textit{#1}}}
\newcommand{\CommentVarTok}[1]{\textcolor[rgb]{0.56,0.35,0.01}{\textbf{\textit{#1}}}}
\newcommand{\ConstantTok}[1]{\textcolor[rgb]{0.56,0.35,0.01}{#1}}
\newcommand{\ControlFlowTok}[1]{\textcolor[rgb]{0.13,0.29,0.53}{\textbf{#1}}}
\newcommand{\DataTypeTok}[1]{\textcolor[rgb]{0.13,0.29,0.53}{#1}}
\newcommand{\DecValTok}[1]{\textcolor[rgb]{0.00,0.00,0.81}{#1}}
\newcommand{\DocumentationTok}[1]{\textcolor[rgb]{0.56,0.35,0.01}{\textbf{\textit{#1}}}}
\newcommand{\ErrorTok}[1]{\textcolor[rgb]{0.64,0.00,0.00}{\textbf{#1}}}
\newcommand{\ExtensionTok}[1]{#1}
\newcommand{\FloatTok}[1]{\textcolor[rgb]{0.00,0.00,0.81}{#1}}
\newcommand{\FunctionTok}[1]{\textcolor[rgb]{0.13,0.29,0.53}{\textbf{#1}}}
\newcommand{\ImportTok}[1]{#1}
\newcommand{\InformationTok}[1]{\textcolor[rgb]{0.56,0.35,0.01}{\textbf{\textit{#1}}}}
\newcommand{\KeywordTok}[1]{\textcolor[rgb]{0.13,0.29,0.53}{\textbf{#1}}}
\newcommand{\NormalTok}[1]{#1}
\newcommand{\OperatorTok}[1]{\textcolor[rgb]{0.81,0.36,0.00}{\textbf{#1}}}
\newcommand{\OtherTok}[1]{\textcolor[rgb]{0.56,0.35,0.01}{#1}}
\newcommand{\PreprocessorTok}[1]{\textcolor[rgb]{0.56,0.35,0.01}{\textit{#1}}}
\newcommand{\RegionMarkerTok}[1]{#1}
\newcommand{\SpecialCharTok}[1]{\textcolor[rgb]{0.81,0.36,0.00}{\textbf{#1}}}
\newcommand{\SpecialStringTok}[1]{\textcolor[rgb]{0.31,0.60,0.02}{#1}}
\newcommand{\StringTok}[1]{\textcolor[rgb]{0.31,0.60,0.02}{#1}}
\newcommand{\VariableTok}[1]{\textcolor[rgb]{0.00,0.00,0.00}{#1}}
\newcommand{\VerbatimStringTok}[1]{\textcolor[rgb]{0.31,0.60,0.02}{#1}}
\newcommand{\WarningTok}[1]{\textcolor[rgb]{0.56,0.35,0.01}{\textbf{\textit{#1}}}}
\usepackage{graphicx}
\makeatletter
\newsavebox\pandoc@box
\newcommand*\pandocbounded[1]{% scales image to fit in text height/width
  \sbox\pandoc@box{#1}%
  \Gscale@div\@tempa{\textheight}{\dimexpr\ht\pandoc@box+\dp\pandoc@box\relax}%
  \Gscale@div\@tempb{\linewidth}{\wd\pandoc@box}%
  \ifdim\@tempb\p@<\@tempa\p@\let\@tempa\@tempb\fi% select the smaller of both
  \ifdim\@tempa\p@<\p@\scalebox{\@tempa}{\usebox\pandoc@box}%
  \else\usebox{\pandoc@box}%
  \fi%
}
% Set default figure placement to htbp
\def\fps@figure{htbp}
\makeatother
\setlength{\emergencystretch}{3em} % prevent overfull lines
\providecommand{\tightlist}{%
  \setlength{\itemsep}{0pt}\setlength{\parskip}{0pt}}
\setcounter{secnumdepth}{-\maxdimen} % remove section numbering
\usepackage{graphicx}
\usepackage{float}
\usepackage{bookmark}
\IfFileExists{xurl.sty}{\usepackage{xurl}}{} % add URL line breaks if available
\urlstyle{same}
\hypersetup{
  hidelinks,
  pdfcreator={LaTeX via pandoc}}

\author{}
\date{\vspace{-2.5em}}

\begin{document}

\begin{Shaded}
\begin{Highlighting}[]
\FunctionTok{options}\NormalTok{(}\AttributeTok{OutDec =} \StringTok{"."}\NormalTok{)  }\CommentTok{\# Asegurar punto decimal en este chunk}

\CommentTok{\# Establecer semilla aleatoria}
\FunctionTok{set.seed}\NormalTok{(}\FunctionTok{sample}\NormalTok{(}\DecValTok{1}\SpecialCharTok{:}\DecValTok{10000}\NormalTok{, }\DecValTok{1}\NormalTok{))}

\CommentTok{\# Aleatorización del contexto del problema}
\NormalTok{contextos }\OtherTok{\textless{}{-}} \FunctionTok{c}\NormalTok{(}
  \StringTok{"empresa"}\NormalTok{, }\StringTok{"organización"}\NormalTok{, }\StringTok{"compañía"}\NormalTok{, }\StringTok{"institución"}\NormalTok{, }\StringTok{"corporación"}\NormalTok{,}
  \StringTok{"entidad"}\NormalTok{, }\StringTok{"firma"}\NormalTok{, }\StringTok{"negocio"}\NormalTok{, }\StringTok{"sociedad"}\NormalTok{, }\StringTok{"grupo empresarial"}
\NormalTok{)}
\NormalTok{contexto }\OtherTok{\textless{}{-}} \FunctionTok{sample}\NormalTok{(contextos, }\DecValTok{1}\NormalTok{)}

\CommentTok{\# Aleatorización de los tipos de bienes}
\NormalTok{tipos\_bien1 }\OtherTok{\textless{}{-}} \FunctionTok{c}\NormalTok{(}\StringTok{"carro"}\NormalTok{, }\StringTok{"auto"}\NormalTok{, }\StringTok{"vehículo"}\NormalTok{, }\StringTok{"automóvil"}\NormalTok{, }\StringTok{"coche"}\NormalTok{)}
\NormalTok{tipos\_bien2 }\OtherTok{\textless{}{-}} \FunctionTok{c}\NormalTok{(}\StringTok{"casa"}\NormalTok{, }\StringTok{"vivienda"}\NormalTok{, }\StringTok{"residencia"}\NormalTok{, }\StringTok{"hogar"}\NormalTok{, }\StringTok{"domicilio"}\NormalTok{)}
\NormalTok{tipos\_bien3 }\OtherTok{\textless{}{-}} \FunctionTok{c}\NormalTok{(}\StringTok{"apartamento"}\NormalTok{, }\StringTok{"departamento"}\NormalTok{, }\StringTok{"piso"}\NormalTok{, }\StringTok{"flat"}\NormalTok{, }\StringTok{"condominio"}\NormalTok{)}

\NormalTok{bien1 }\OtherTok{\textless{}{-}} \FunctionTok{sample}\NormalTok{(tipos\_bien1, }\DecValTok{1}\NormalTok{)}
\NormalTok{bien2 }\OtherTok{\textless{}{-}} \FunctionTok{sample}\NormalTok{(tipos\_bien2, }\DecValTok{1}\NormalTok{)}
\NormalTok{bien3 }\OtherTok{\textless{}{-}} \FunctionTok{sample}\NormalTok{(tipos\_bien3, }\DecValTok{1}\NormalTok{)}

\CommentTok{\# Aleatorización de términos para el enunciado}
\NormalTok{terminos\_encuesta }\OtherTok{\textless{}{-}} \FunctionTok{c}\NormalTok{(}\StringTok{"encuesta"}\NormalTok{, }\StringTok{"sondeo"}\NormalTok{, }\StringTok{"estudio"}\NormalTok{, }\StringTok{"investigación"}\NormalTok{, }\StringTok{"consulta"}\NormalTok{)}
\NormalTok{termino\_encuesta }\OtherTok{\textless{}{-}} \FunctionTok{sample}\NormalTok{(terminos\_encuesta, }\DecValTok{1}\NormalTok{)}

\NormalTok{terminos\_personas }\OtherTok{\textless{}{-}} \FunctionTok{c}\NormalTok{(}\StringTok{"personas"}\NormalTok{, }\StringTok{"empleados"}\NormalTok{, }\StringTok{"trabajadores"}\NormalTok{, }\StringTok{"colaboradores"}\NormalTok{, }\StringTok{"miembros"}\NormalTok{)}
\NormalTok{termino\_personas }\OtherTok{\textless{}{-}} \FunctionTok{sample}\NormalTok{(terminos\_personas, }\DecValTok{1}\NormalTok{)}

\NormalTok{terminos\_bienes }\OtherTok{\textless{}{-}} \FunctionTok{c}\NormalTok{(}\StringTok{"bienes"}\NormalTok{, }\StringTok{"posesiones"}\NormalTok{, }\StringTok{"propiedades"}\NormalTok{, }\StringTok{"activos"}\NormalTok{, }\StringTok{"patrimonios"}\NormalTok{)}
\NormalTok{termino\_bienes }\OtherTok{\textless{}{-}} \FunctionTok{sample}\NormalTok{(terminos\_bienes, }\DecValTok{1}\NormalTok{)}

\NormalTok{terminos\_resultados }\OtherTok{\textless{}{-}} \FunctionTok{c}\NormalTok{(}\StringTok{"resultados"}\NormalTok{, }\StringTok{"datos"}\NormalTok{, }\StringTok{"información"}\NormalTok{, }\StringTok{"estadísticas"}\NormalTok{, }\StringTok{"cifras"}\NormalTok{)}
\NormalTok{termino\_resultados }\OtherTok{\textless{}{-}} \FunctionTok{sample}\NormalTok{(terminos\_resultados, }\DecValTok{1}\NormalTok{)}

\CommentTok{\# Aleatorización de colores para el gráfico circular}
\NormalTok{paleta\_colores }\OtherTok{\textless{}{-}} \FunctionTok{list}\NormalTok{(}
  \FunctionTok{c}\NormalTok{(}\StringTok{"\#E57373"}\NormalTok{, }\StringTok{"\#81C784"}\NormalTok{, }\StringTok{"\#64B5F6"}\NormalTok{, }\StringTok{"\#FFB74D"}\NormalTok{, }\StringTok{"\#BA68C8"}\NormalTok{),  }\CommentTok{\# Paleta pastel}
  \FunctionTok{c}\NormalTok{(}\StringTok{"\#D32F2F"}\NormalTok{, }\StringTok{"\#388E3C"}\NormalTok{, }\StringTok{"\#1976D2"}\NormalTok{, }\StringTok{"\#F57C00"}\NormalTok{, }\StringTok{"\#7B1FA2"}\NormalTok{),  }\CommentTok{\# Paleta saturada}
  \FunctionTok{c}\NormalTok{(}\StringTok{"\#FFCDD2"}\NormalTok{, }\StringTok{"\#C8E6C9"}\NormalTok{, }\StringTok{"\#BBDEFB"}\NormalTok{, }\StringTok{"\#FFE0B2"}\NormalTok{, }\StringTok{"\#E1BEE7"}\NormalTok{),  }\CommentTok{\# Paleta suave}
  \FunctionTok{c}\NormalTok{(}\StringTok{"\#EF9A9A"}\NormalTok{, }\StringTok{"\#A5D6A7"}\NormalTok{, }\StringTok{"\#90CAF9"}\NormalTok{, }\StringTok{"\#FFCC80"}\NormalTok{, }\StringTok{"\#CE93D8"}\NormalTok{),  }\CommentTok{\# Paleta media}
  \FunctionTok{c}\NormalTok{(}\StringTok{"\#B71C1C"}\NormalTok{, }\StringTok{"\#1B5E20"}\NormalTok{, }\StringTok{"\#0D47A1"}\NormalTok{, }\StringTok{"\#E65100"}\NormalTok{, }\StringTok{"\#4A148C"}\NormalTok{)   }\CommentTok{\# Paleta oscura}
\NormalTok{)}
\NormalTok{paleta\_seleccionada }\OtherTok{\textless{}{-}} \FunctionTok{sample}\NormalTok{(paleta\_colores, }\DecValTok{1}\NormalTok{)[[}\DecValTok{1}\NormalTok{]]}

\CommentTok{\# Aleatorización de los porcentajes manteniendo coherencia matemática}
\CommentTok{\# Generación de porcentajes iniciales para cada categoría}
\NormalTok{set\_porcentajes }\OtherTok{\textless{}{-}} \ControlFlowTok{function}\NormalTok{() \{}
  \ControlFlowTok{repeat}\NormalTok{ \{}
    \CommentTok{\# Generar valores base aleatorios}
\NormalTok{    p\_bien1\_bien3 }\OtherTok{\textless{}{-}} \FunctionTok{sample}\NormalTok{(}\DecValTok{15}\SpecialCharTok{:}\DecValTok{35}\NormalTok{, }\DecValTok{1}\NormalTok{)  }\CommentTok{\# Porcentaje de bien1 y bien3}
\NormalTok{    p\_solo\_bien3 }\OtherTok{\textless{}{-}} \FunctionTok{sample}\NormalTok{(}\DecValTok{15}\SpecialCharTok{:}\DecValTok{30}\NormalTok{, }\DecValTok{1}\NormalTok{)   }\CommentTok{\# Porcentaje solo bien3}
\NormalTok{    p\_solo\_bien2 }\OtherTok{\textless{}{-}} \FunctionTok{sample}\NormalTok{(}\DecValTok{25}\SpecialCharTok{:}\DecValTok{40}\NormalTok{, }\DecValTok{1}\NormalTok{)   }\CommentTok{\# Porcentaje solo bien2}
\NormalTok{    p\_solo\_bien1 }\OtherTok{\textless{}{-}} \FunctionTok{sample}\NormalTok{(}\DecValTok{5}\SpecialCharTok{:}\DecValTok{10}\NormalTok{, }\DecValTok{1}\NormalTok{)    }\CommentTok{\# Porcentaje solo bien1}

    \CommentTok{\# Calcular el porcentaje restante para bien1 y bien2}
\NormalTok{    p\_bien1\_bien2 }\OtherTok{\textless{}{-}} \DecValTok{100} \SpecialCharTok{{-}}\NormalTok{ (p\_bien1\_bien3 }\SpecialCharTok{+}\NormalTok{ p\_solo\_bien3 }\SpecialCharTok{+}\NormalTok{ p\_solo\_bien2 }\SpecialCharTok{+}\NormalTok{ p\_solo\_bien1)}

    \CommentTok{\# Validar que el porcentaje restante sea positivo y razonable}
    \ControlFlowTok{if}\NormalTok{ (p\_bien1\_bien2 }\SpecialCharTok{\textgreater{}=} \DecValTok{10} \SpecialCharTok{\&\&}\NormalTok{ p\_bien1\_bien2 }\SpecialCharTok{\textless{}=} \DecValTok{25}\NormalTok{) \{}
      \FunctionTok{return}\NormalTok{(}\FunctionTok{c}\NormalTok{(p\_bien1\_bien3, p\_solo\_bien3, p\_solo\_bien2, p\_solo\_bien1, p\_bien1\_bien2))}
\NormalTok{    \}}
\NormalTok{  \}}
\NormalTok{\}}

\CommentTok{\# Obtener porcentajes válidos}
\NormalTok{porcentajes }\OtherTok{\textless{}{-}} \FunctionTok{set\_porcentajes}\NormalTok{()}

\NormalTok{p\_bien1\_bien3 }\OtherTok{\textless{}{-}}\NormalTok{ porcentajes[}\DecValTok{1}\NormalTok{]  }\CommentTok{\# Porcentaje de bien1 y bien3}
\NormalTok{p\_solo\_bien3 }\OtherTok{\textless{}{-}}\NormalTok{ porcentajes[}\DecValTok{2}\NormalTok{]   }\CommentTok{\# Porcentaje solo bien3}
\NormalTok{p\_solo\_bien2 }\OtherTok{\textless{}{-}}\NormalTok{ porcentajes[}\DecValTok{3}\NormalTok{]   }\CommentTok{\# Porcentaje solo bien2}
\NormalTok{p\_solo\_bien1 }\OtherTok{\textless{}{-}}\NormalTok{ porcentajes[}\DecValTok{4}\NormalTok{]   }\CommentTok{\# Porcentaje solo bien1}
\NormalTok{p\_bien1\_bien2 }\OtherTok{\textless{}{-}}\NormalTok{ porcentajes[}\DecValTok{5}\NormalTok{]  }\CommentTok{\# Porcentaje de bien1 y bien2}

\CommentTok{\# Verificar que suman 100\%}
\FunctionTok{test\_that}\NormalTok{(}\StringTok{"Los porcentajes suman 100\%"}\NormalTok{, \{}
  \FunctionTok{expect\_equal}\NormalTok{(}\FunctionTok{sum}\NormalTok{(porcentajes), }\DecValTok{100}\NormalTok{)}
\NormalTok{\})}
\end{Highlighting}
\end{Shaded}

Test passed

\begin{Shaded}
\begin{Highlighting}[]
\CommentTok{\# Valor conocido para el cálculo: Número de personas con bien1 y bien3}
\NormalTok{personas\_bien1\_bien3 }\OtherTok{\textless{}{-}} \FunctionTok{sample}\NormalTok{(}\FunctionTok{c}\NormalTok{(}\DecValTok{72}\NormalTok{, }\DecValTok{96}\NormalTok{, }\DecValTok{120}\NormalTok{, }\DecValTok{144}\NormalTok{, }\DecValTok{168}\NormalTok{, }\DecValTok{192}\NormalTok{, }\DecValTok{216}\NormalTok{, }\DecValTok{240}\NormalTok{), }\DecValTok{1}\NormalTok{)}

\CommentTok{\# Calcular el total de personas}
\NormalTok{total\_personas }\OtherTok{\textless{}{-}} \FunctionTok{round}\NormalTok{(personas\_bien1\_bien3 }\SpecialCharTok{*} \DecValTok{100} \SpecialCharTok{/}\NormalTok{ p\_bien1\_bien3)}

\CommentTok{\# Calcular el número de personas en cada categoría}
\NormalTok{personas\_solo\_bien3 }\OtherTok{\textless{}{-}} \FunctionTok{round}\NormalTok{(total\_personas }\SpecialCharTok{*}\NormalTok{ p\_solo\_bien3 }\SpecialCharTok{/} \DecValTok{100}\NormalTok{)}
\NormalTok{personas\_solo\_bien2 }\OtherTok{\textless{}{-}} \FunctionTok{round}\NormalTok{(total\_personas }\SpecialCharTok{*}\NormalTok{ p\_solo\_bien2 }\SpecialCharTok{/} \DecValTok{100}\NormalTok{)}
\NormalTok{personas\_solo\_bien1 }\OtherTok{\textless{}{-}} \FunctionTok{round}\NormalTok{(total\_personas }\SpecialCharTok{*}\NormalTok{ p\_solo\_bien1 }\SpecialCharTok{/} \DecValTok{100}\NormalTok{)}
\NormalTok{personas\_bien1\_bien2 }\OtherTok{\textless{}{-}} \FunctionTok{round}\NormalTok{(total\_personas }\SpecialCharTok{*}\NormalTok{ p\_bien1\_bien2 }\SpecialCharTok{/} \DecValTok{100}\NormalTok{)}

\CommentTok{\# Recalcular personas\_bien1\_bien3 para asegurar que el total sea correcto}
\NormalTok{personas\_bien1\_bien3 }\OtherTok{\textless{}{-}}\NormalTok{ total\_personas }\SpecialCharTok{{-}}\NormalTok{ (personas\_solo\_bien3 }\SpecialCharTok{+}\NormalTok{ personas\_solo\_bien2 }\SpecialCharTok{+}\NormalTok{ personas\_solo\_bien1 }\SpecialCharTok{+}\NormalTok{ personas\_bien1\_bien2)}

\CommentTok{\# Asegurar que todos los valores son positivos y razonables}
\NormalTok{personas }\OtherTok{\textless{}{-}} \FunctionTok{c}\NormalTok{(personas\_bien1\_bien3, personas\_solo\_bien3, personas\_solo\_bien2, personas\_solo\_bien1, personas\_bien1\_bien2)}
\FunctionTok{test\_that}\NormalTok{(}\StringTok{"Todas las categorías tienen al menos una persona"}\NormalTok{, \{}
  \FunctionTok{expect\_true}\NormalTok{(}\FunctionTok{all}\NormalTok{(personas }\SpecialCharTok{\textgreater{}} \DecValTok{0}\NormalTok{))}
\NormalTok{\})}
\end{Highlighting}
\end{Shaded}

Test passed

\begin{Shaded}
\begin{Highlighting}[]
\CommentTok{\# Asegurar que la suma de todas las categorías es igual al total}
\FunctionTok{test\_that}\NormalTok{(}\StringTok{"La suma de todas las categorías es igual al total"}\NormalTok{, \{}
  \FunctionTok{expect\_equal}\NormalTok{(}\FunctionTok{sum}\NormalTok{(personas), total\_personas)}
\NormalTok{\})}
\end{Highlighting}
\end{Shaded}

Test passed

\begin{Shaded}
\begin{Highlighting}[]
\CommentTok{\# Recalcular los porcentajes reales basados en los números de personas (para mantener coherencia)}
\NormalTok{p\_bien1\_bien3 }\OtherTok{\textless{}{-}} \FunctionTok{round}\NormalTok{(personas\_bien1\_bien3 }\SpecialCharTok{/}\NormalTok{ total\_personas }\SpecialCharTok{*} \DecValTok{100}\NormalTok{)}
\NormalTok{p\_solo\_bien3 }\OtherTok{\textless{}{-}} \FunctionTok{round}\NormalTok{(personas\_solo\_bien3 }\SpecialCharTok{/}\NormalTok{ total\_personas }\SpecialCharTok{*} \DecValTok{100}\NormalTok{)}
\NormalTok{p\_solo\_bien2 }\OtherTok{\textless{}{-}} \FunctionTok{round}\NormalTok{(personas\_solo\_bien2 }\SpecialCharTok{/}\NormalTok{ total\_personas }\SpecialCharTok{*} \DecValTok{100}\NormalTok{)}
\NormalTok{p\_solo\_bien1 }\OtherTok{\textless{}{-}} \FunctionTok{round}\NormalTok{(personas\_solo\_bien1 }\SpecialCharTok{/}\NormalTok{ total\_personas }\SpecialCharTok{*} \DecValTok{100}\NormalTok{)}
\NormalTok{p\_bien1\_bien2 }\OtherTok{\textless{}{-}} \FunctionTok{round}\NormalTok{(personas\_bien1\_bien2 }\SpecialCharTok{/}\NormalTok{ total\_personas }\SpecialCharTok{*} \DecValTok{100}\NormalTok{)}

\CommentTok{\# Ajustar para asegurar que suman 100\%}
\NormalTok{total\_porcentaje }\OtherTok{\textless{}{-}}\NormalTok{ p\_bien1\_bien3 }\SpecialCharTok{+}\NormalTok{ p\_solo\_bien3 }\SpecialCharTok{+}\NormalTok{ p\_solo\_bien2 }\SpecialCharTok{+}\NormalTok{ p\_solo\_bien1 }\SpecialCharTok{+}\NormalTok{ p\_bien1\_bien2}
\ControlFlowTok{if}\NormalTok{ (total\_porcentaje }\SpecialCharTok{\textgreater{}} \DecValTok{100}\NormalTok{) \{}
  \CommentTok{\# Si suman más de 100, restar la diferencia de la categoría más grande}
\NormalTok{  max\_idx }\OtherTok{\textless{}{-}} \FunctionTok{which.max}\NormalTok{(porcentajes)}
\NormalTok{  porcentajes[max\_idx] }\OtherTok{\textless{}{-}}\NormalTok{ porcentajes[max\_idx] }\SpecialCharTok{{-}}\NormalTok{ (total\_porcentaje }\SpecialCharTok{{-}} \DecValTok{100}\NormalTok{)}
\NormalTok{\} }\ControlFlowTok{else} \ControlFlowTok{if}\NormalTok{ (total\_porcentaje }\SpecialCharTok{\textless{}} \DecValTok{100}\NormalTok{) \{}
  \CommentTok{\# Si suman menos de 100, añadir la diferencia a la categoría más pequeña}
\NormalTok{  min\_idx }\OtherTok{\textless{}{-}} \FunctionTok{which.min}\NormalTok{(porcentajes)}
\NormalTok{  porcentajes[min\_idx] }\OtherTok{\textless{}{-}}\NormalTok{ porcentajes[min\_idx] }\SpecialCharTok{+}\NormalTok{ (}\DecValTok{100} \SpecialCharTok{{-}}\NormalTok{ total\_porcentaje)}
\NormalTok{\}}

\NormalTok{p\_bien1\_bien3 }\OtherTok{\textless{}{-}}\NormalTok{ porcentajes[}\DecValTok{1}\NormalTok{]}
\NormalTok{p\_solo\_bien3 }\OtherTok{\textless{}{-}}\NormalTok{ porcentajes[}\DecValTok{2}\NormalTok{]}
\NormalTok{p\_solo\_bien2 }\OtherTok{\textless{}{-}}\NormalTok{ porcentajes[}\DecValTok{3}\NormalTok{]}
\NormalTok{p\_solo\_bien1 }\OtherTok{\textless{}{-}}\NormalTok{ porcentajes[}\DecValTok{4}\NormalTok{]}
\NormalTok{p\_bien1\_bien2 }\OtherTok{\textless{}{-}}\NormalTok{ porcentajes[}\DecValTok{5}\NormalTok{]}

\CommentTok{\# Asegurar que los porcentajes suman exactamente 100\%}
\FunctionTok{test\_that}\NormalTok{(}\StringTok{"Los porcentajes ajustados suman 100\%"}\NormalTok{, \{}
  \FunctionTok{expect\_equal}\NormalTok{(}\FunctionTok{sum}\NormalTok{(porcentajes), }\DecValTok{100}\NormalTok{)}
\NormalTok{\})}
\end{Highlighting}
\end{Shaded}

Test passed

\begin{Shaded}
\begin{Highlighting}[]
\CommentTok{\# Determinar la respuesta correcta y generar tres distractores plausibles}
\CommentTok{\# La respuesta correcta es el número de personas con solo bien3}
\NormalTok{respuesta\_correcta }\OtherTok{\textless{}{-}}\NormalTok{ personas\_solo\_bien3}

\CommentTok{\# Generar distractores plausibles}
\NormalTok{distractor1 }\OtherTok{\textless{}{-}} \FunctionTok{round}\NormalTok{(personas\_bien1\_bien3 }\SpecialCharTok{/} \DecValTok{4}\NormalTok{)  }\CommentTok{\# Un cuarto del valor dado en el problema}
\NormalTok{distractor2 }\OtherTok{\textless{}{-}} \FunctionTok{round}\NormalTok{(total\_personas }\SpecialCharTok{*}\NormalTok{ p\_bien1\_bien3 }\SpecialCharTok{/} \DecValTok{100}\NormalTok{)  }\CommentTok{\# Confundir entre porcentaje y valor absoluto}
\NormalTok{distractor3 }\OtherTok{\textless{}{-}}\NormalTok{ personas\_bien1\_bien3  }\CommentTok{\# El mismo valor dado en el problema}

\CommentTok{\# Asegurarse de que todos los distractores son diferentes de la respuesta correcta}
\ControlFlowTok{if}\NormalTok{ (distractor1 }\SpecialCharTok{==}\NormalTok{ respuesta\_correcta) distractor1 }\OtherTok{\textless{}{-}}\NormalTok{ distractor1 }\SpecialCharTok{+} \FunctionTok{sample}\NormalTok{(}\DecValTok{5}\SpecialCharTok{:}\DecValTok{15}\NormalTok{, }\DecValTok{1}\NormalTok{)}
\ControlFlowTok{if}\NormalTok{ (distractor2 }\SpecialCharTok{==}\NormalTok{ respuesta\_correcta) distractor2 }\OtherTok{\textless{}{-}}\NormalTok{ distractor2 }\SpecialCharTok{{-}} \FunctionTok{sample}\NormalTok{(}\DecValTok{5}\SpecialCharTok{:}\DecValTok{15}\NormalTok{, }\DecValTok{1}\NormalTok{)}
\ControlFlowTok{if}\NormalTok{ (distractor3 }\SpecialCharTok{==}\NormalTok{ respuesta\_correcta) distractor3 }\OtherTok{\textless{}{-}}\NormalTok{ distractor3 }\SpecialCharTok{+} \FunctionTok{sample}\NormalTok{(}\DecValTok{5}\SpecialCharTok{:}\DecValTok{15}\NormalTok{, }\DecValTok{1}\NormalTok{)}

\CommentTok{\# Asegurarse de que todos los distractores son diferentes entre sí}
\ControlFlowTok{while}\NormalTok{ (}\FunctionTok{length}\NormalTok{(}\FunctionTok{unique}\NormalTok{(}\FunctionTok{c}\NormalTok{(distractor1, distractor2, distractor3))) }\SpecialCharTok{\textless{}} \DecValTok{3}\NormalTok{) \{}
  \ControlFlowTok{if}\NormalTok{ (distractor1 }\SpecialCharTok{==}\NormalTok{ distractor2) distractor1 }\OtherTok{\textless{}{-}}\NormalTok{ distractor1 }\SpecialCharTok{+} \FunctionTok{sample}\NormalTok{(}\DecValTok{5}\SpecialCharTok{:}\DecValTok{10}\NormalTok{, }\DecValTok{1}\NormalTok{)}
  \ControlFlowTok{if}\NormalTok{ (distractor2 }\SpecialCharTok{==}\NormalTok{ distractor3) distractor2 }\OtherTok{\textless{}{-}}\NormalTok{ distractor2 }\SpecialCharTok{{-}} \FunctionTok{sample}\NormalTok{(}\DecValTok{5}\SpecialCharTok{:}\DecValTok{10}\NormalTok{, }\DecValTok{1}\NormalTok{)}
  \ControlFlowTok{if}\NormalTok{ (distractor1 }\SpecialCharTok{==}\NormalTok{ distractor3) distractor3 }\OtherTok{\textless{}{-}}\NormalTok{ distractor3 }\SpecialCharTok{+} \FunctionTok{sample}\NormalTok{(}\DecValTok{5}\SpecialCharTok{:}\DecValTok{10}\NormalTok{, }\DecValTok{1}\NormalTok{)}
\NormalTok{\}}

\CommentTok{\# Crear un vector con todas las opciones y mezclarlas}
\NormalTok{opciones }\OtherTok{\textless{}{-}} \FunctionTok{c}\NormalTok{(respuesta\_correcta, distractor1, distractor2, distractor3)}
\FunctionTok{names}\NormalTok{(opciones) }\OtherTok{\textless{}{-}} \FunctionTok{c}\NormalTok{(}\StringTok{"correcta"}\NormalTok{, }\StringTok{"distractor1"}\NormalTok{, }\StringTok{"distractor2"}\NormalTok{, }\StringTok{"distractor3"}\NormalTok{)}
\NormalTok{opciones\_mezcladas }\OtherTok{\textless{}{-}} \FunctionTok{sample}\NormalTok{(opciones)}

\CommentTok{\# Identificar la posición de la respuesta correcta en las opciones mezcladas}
\NormalTok{indice\_correcto }\OtherTok{\textless{}{-}} \FunctionTok{which}\NormalTok{(opciones\_mezcladas }\SpecialCharTok{==}\NormalTok{ respuesta\_correcta)}

\CommentTok{\# Crear el vector de solución para r{-}exams}
\NormalTok{solucion }\OtherTok{\textless{}{-}} \FunctionTok{rep}\NormalTok{(}\DecValTok{0}\NormalTok{, }\DecValTok{4}\NormalTok{)}
\NormalTok{solucion[indice\_correcto] }\OtherTok{\textless{}{-}} \DecValTok{1}
\end{Highlighting}
\end{Shaded}

\begin{Shaded}
\begin{Highlighting}[]
\FunctionTok{options}\NormalTok{(}\AttributeTok{OutDec =} \StringTok{"."}\NormalTok{)  }\CommentTok{\# Asegurar punto decimal en este chunk}

\CommentTok{\# Código Python para generar el gráfico circular}
\NormalTok{codigo\_python }\OtherTok{\textless{}{-}} \FunctionTok{paste0}\NormalTok{(}\StringTok{"}
\StringTok{import matplotlib}
\StringTok{matplotlib.use(\textquotesingle{}Agg\textquotesingle{})  \# Usar backend no interactivo}
\StringTok{import matplotlib.pyplot as plt}
\StringTok{import numpy as np}

\StringTok{\# Datos para el gráfico}
\StringTok{labels = [\textquotesingle{}"}\NormalTok{, bien1, }\StringTok{" y "}\NormalTok{, bien3, }\StringTok{"\textquotesingle{}, \textquotesingle{}Solo "}\NormalTok{, bien3, }\StringTok{"\textquotesingle{}, \textquotesingle{}Solo "}\NormalTok{, bien2, }\StringTok{"\textquotesingle{},}
\StringTok{          \textquotesingle{}Solo "}\NormalTok{, bien1, }\StringTok{"\textquotesingle{}, \textquotesingle{}"}\NormalTok{, bien1, }\StringTok{" y "}\NormalTok{, bien2, }\StringTok{"\textquotesingle{}]}
\StringTok{sizes = ["}\NormalTok{, p\_bien1\_bien3, }\StringTok{", "}\NormalTok{, p\_solo\_bien3, }\StringTok{", "}\NormalTok{, p\_solo\_bien2, }\StringTok{",}
\StringTok{         "}\NormalTok{, p\_solo\_bien1, }\StringTok{", "}\NormalTok{, p\_bien1\_bien2, }\StringTok{"]}
\StringTok{colors = [\textquotesingle{}"}\NormalTok{, paleta\_seleccionada[}\DecValTok{1}\NormalTok{], }\StringTok{"\textquotesingle{}, \textquotesingle{}"}\NormalTok{, paleta\_seleccionada[}\DecValTok{2}\NormalTok{], }\StringTok{"\textquotesingle{},}
\StringTok{          \textquotesingle{}"}\NormalTok{, paleta\_seleccionada[}\DecValTok{3}\NormalTok{], }\StringTok{"\textquotesingle{}, \textquotesingle{}"}\NormalTok{, paleta\_seleccionada[}\DecValTok{4}\NormalTok{], }\StringTok{"\textquotesingle{},}
\StringTok{          \textquotesingle{}"}\NormalTok{, paleta\_seleccionada[}\DecValTok{5}\NormalTok{], }\StringTok{"\textquotesingle{}]}

\StringTok{\# Explode para destacar ligeramente el sector con bien1 y bien3}
\StringTok{explode = (0.03, 0, 0, 0, 0)}

\StringTok{\# Crear figura con tamaño controlado}
\StringTok{plt.figure(figsize=(5, 4))}

\StringTok{\# Crear gráfico circular con bordes blancos}
\StringTok{patches, texts, autotexts = plt.pie(}
\StringTok{    sizes,}
\StringTok{    explode=explode,}
\StringTok{    colors=colors,}
\StringTok{    shadow=True,}
\StringTok{    startangle=90,}
\StringTok{    autopct=\textquotesingle{}\%d\%\%\textquotesingle{},}
\StringTok{    pctdistance=0.85,}
\StringTok{    wedgeprops=\{\textquotesingle{}edgecolor\textquotesingle{}: \textquotesingle{}white\textquotesingle{}, \textquotesingle{}linewidth\textquotesingle{}: 1\}}
\StringTok{)}

\StringTok{\# Configuración estética de los textos de porcentaje}
\StringTok{for autotext in autotexts:}
\StringTok{    autotext.set\_fontsize(9)}
\StringTok{    autotext.set\_weight(\textquotesingle{}bold\textquotesingle{})}
\StringTok{    autotext.set\_color(\textquotesingle{}white\textquotesingle{})}

\StringTok{\# Crear leyenda por separado}
\StringTok{plt.axis(\textquotesingle{}equal\textquotesingle{})  \# Asegurar que el gráfico sea un círculo}
\StringTok{plt.legend(}
\StringTok{    labels,}
\StringTok{    loc=\textquotesingle{}upper left\textquotesingle{},}
\StringTok{    bbox\_to\_anchor=({-}0.1, 1),}
\StringTok{    fontsize=8,}
\StringTok{    frameon=False}
\StringTok{)}

\StringTok{\# Añadir título centrado en la parte superior}
\StringTok{plt.title(\textquotesingle{}Resultados de la encuesta\textquotesingle{}, fontsize=10)}

\StringTok{plt.tight\_layout()}

\StringTok{\# Guardar en múltiples formatos para asegurar compatibilidad}
\StringTok{plt.savefig(\textquotesingle{}grafico\_circular.png\textquotesingle{}, dpi=150, bbox\_inches=\textquotesingle{}tight\textquotesingle{},}
\StringTok{           transparent=True, format=\textquotesingle{}png\textquotesingle{})}
\StringTok{plt.savefig(\textquotesingle{}grafico\_circular.pdf\textquotesingle{}, dpi=150, bbox\_inches=\textquotesingle{}tight\textquotesingle{},}
\StringTok{           transparent=True, format=\textquotesingle{}pdf\textquotesingle{})}
\StringTok{plt.close()}
\StringTok{"}\NormalTok{)}

\CommentTok{\# Ejecutar código Python para generar la figura}
\FunctionTok{py\_run\_string}\NormalTok{(codigo\_python)}
\end{Highlighting}
\end{Shaded}

\section{Question}\label{question}

Se realizó un(a) investigación a un grupo de empleados de un(a) firma
sobre el tipo de posesiones que poseen. Los estadísticas se presentan en
la gráfica.

\includegraphics[width=0.7\linewidth,height=\textheight,keepaspectratio]{grafico_circular.png}

Si 119 empleados de la firma poseen automóvil y condominio, ¿cuántas
personas poseen solo condominio?

\subsection{Answerlist}\label{answerlist}

\begin{itemize}
\tightlist
\item
  120
\item
  30
\item
  218
\item
  119
\end{itemize}

\section{Solution}\label{solution}

Para resolver este problema, necesitamos aplicar proporciones a partir
de la información dada:

\begin{enumerate}
\def\labelenumi{\arabic{enumi})}
\item
  Sabemos que 119 empleados poseen automóvil y condominio, lo que
  representa el 16\% del total según el gráfico circular.
\item
  Primero calculamos el número total de empleados en la firma:

  Total de empleados = (119 × 100\%) ÷ 16\% = 750
\item
  Según el gráfico, los que poseen solo condominio representan el 29\%
  del total.
\item
  Por lo tanto, el número de empleados que poseen solo condominio es:

  29\% × 750 = 218
\end{enumerate}

Así, 218 empleados de la firma poseen solo condominio.

\subsection{Answerlist}\label{answerlist-1}

\begin{itemize}
\tightlist
\item
  Falso
\item
  Falso
\item
  Verdadero
\item
  Falso
\end{itemize}

\section{Meta-information}\label{meta-information}

exname: proporciones\_diagrama\_circular extype: schoice exsolution:
0010 exshuffle: TRUE exsection:
Estadística\textbar Proporciones\textbar Interpretación de gráficos

\end{document}
